\documentclass[11pt]{article}
\usepackage{amsmath,amsthm,amssymb,amsfonts}
\usepackage[margin=1in]{geometry}
\usepackage[colorlinks=true,linkcolor=blue,citecolor=blue,urlcolor=blue]{hyperref}

\title{Three Classical Analytic Estimates for Explicit-Formula Coefficients of the Riemann Zeta Function}
\author{Robert Duran\thanks{Email: \texttt{duran.robert301@gmail.com}.}}
\date{}

\theoremstyle{plain}
\newtheorem{theorem}{Theorem}[section]
\newtheorem{lemma}[theorem]{Lemma}
\newtheorem{corollary}[theorem]{Corollary}

\theoremstyle{definition}
\newtheorem{definition}[theorem]{Definition}
\newtheorem{remark}[theorem]{Remark}

\begin{document}
\maketitle

\begin{abstract}
We formulate three classical analytic estimates for coefficient families arising in explicit-formula expansions indexed by the non-trivial zeros of the Riemann zeta function: (i) a dyadic $L^2$ tail bound, (ii) a large-sieve--type mean-square inequality for exponential sums over zero ordinates \emph{without} any global spacing hypothesis, and (iii) a uniform mean-square tail bound for a weighted explicit-formula expansion. The arguments rely only on Stirling asymptotics for $\Gamma$, the Riemann--von Mangoldt formula (including unit-interval zero counts), and a classical large-sieve inequality for separated exponentials.
\end{abstract}

\section{Introduction}
Let $\rho = \beta + i\gamma$ range over the non-trivial zeros of the Riemann zeta function, paired under the functional equation. Write the explicit-formula contribution to $\psi(x)$ in the form
\[
  \log R(s) = \sum_{\rho} a_\rho\, e^{(\rho-1)s} + \cdots, \qquad
  j_{\mathrm{rel}}(s) = \sum_{\rho} \frac{\rho-1}{\rho}\, a_\rho\, e^{(\rho-1)s} + \cdots,
\]
and define coefficients
\[
  b_\rho := \frac{\rho-1}{\rho}\, G(\rho), \qquad
  G(s) := \pi^{-s/2} \Gamma\!\left(\frac{s}{2}\right) H(s).
\]
Here $H(s)$ is a holomorphic factor coming from Mellin transforms of test functions, rational weights, and possible $\zeta'/\zeta$ terms, assumed to be of at most polynomial growth in vertical strips with $\Re(s+\kappa)\ge 1/2+\varepsilon$ for some $\kappa,\varepsilon>0$. We present three self-contained analytic estimates for these coefficients and their dyadic aggregates. No reference is made to ERURH or to RH in what follows.

\section{Standing hypotheses and preliminaries}
Let $\rho=\beta+i\gamma$ range over the non-trivial zeros of $\zeta(s)$, counted with multiplicity, and let
\[
  b_\rho := \frac{\rho-1}{\rho}\,G(\rho),
  \qquad
  G(s) := \pi^{-s/2}\Gamma\!\left(\frac{s}{2}\right)H(s).
\]
The estimates in this note are stated and proved under the following \emph{explicit} hypotheses, which isolate exactly the analytic input required.

\subsection*{Standing hypotheses}
\noindent\textbf{(H1) (Holomorphic model with gamma factor).}
The auxiliary factor $H(s)$ is holomorphic in a neighbourhood of the closed strip $0\le \Re s\le 1$, and the coefficients $b_\rho$ are defined by the above formula (in particular, the completed zeta-factor $\pi^{-s/2}\Gamma(s/2)$ is present in $G(s)$).

\medskip
\noindent\textbf{(H2) (Polynomial vertical growth of $H$).}
There exist constants $C_H>0$ and $A_H\ge 0$ such that
\[
  |H(\sigma+it)| \le C_H(1+|t|)^{A_H}
  \qquad (0\le \sigma\le 1,\ t\in\mathbb{R}).
\]

\medskip
\noindent\textbf{(H3) (Riemann--von Mangoldt and unit-interval zero counts).}
Let $N(T)$ denote the number of non-trivial zeros $\rho=\beta+i\gamma$ with $0<\beta<1$ and $0<\gamma\le T$, counted with multiplicity. We use the classical Riemann--von Mangoldt formula and its consequence that, for $|y|\ge 3$,
\[
  \#\{\rho:\ y<\Im(\rho)\le y+1\}\ \ll\ \log(|y|+2),
\]
with an absolute implied constant.

\begin{remark}[Verifying the hypotheses in explicit-formula settings]\label{rem:verify-hyp}
Hypothesis \textbf{(H1)} is automatic once one works with the completed zeta function
$\Lambda(s)=\pi^{-s/2}\Gamma(s/2)\zeta(s)$, from which the gamma factor enters the explicit formula.
Hypothesis \textbf{(H2)} holds (with room to spare) when $H(s)$ is a Mellin transform of a smooth compactly supported test function, since such Mellin transforms are entire and rapidly decaying in vertical strips. Rational weights contribute at most polynomial growth.

If one wishes to incorporate additional factors such as $(\zeta'/\zeta)(s+\kappa)$ into $H(s)$, then for the present arguments it suffices that the resulting $H$ remain holomorphic in a neighbourhood of $0\le \Re s\le 1$ and satisfy \textbf{(H2)} there. A simple sufficient condition is $\kappa>1$, which places $s+\kappa$ in the zero-free half-plane $\Re>1$ for all $s$ in the strip.
\end{remark}

\medskip
All implied constants may depend on $C_H,A_H$ (and hence on the specific choice of $H$), but are otherwise absolute.

\medskip
We record three standard lemmas (Stirling decay, local zero counting, and a classical large-sieve inequality for separated exponentials) that will be used repeatedly.
\medskip

\begin{lemma}[Stirling decay for the coefficients]\label{lem:decay}
Assume Hypotheses \textbf{(H1)}--\textbf{(H2)}. In particular, $H(s)$ is holomorphic in a neighbourhood of the strip $0\le \Re s\le 1$ and, for some $A_H\ge 0$,
\[
  H(\sigma+it) \ll (1+|t|)^{A_H}
  \qquad (0\le \sigma\le 1,\ t\in\mathbb{R}).
\]
Then there exist constants $c>0$ and $k\ge 0$ (depending only on $A_H$) such that for every non-trivial zero $\rho=\beta+i\gamma$ of $\zeta(s)$,
\[
  |b_\rho| \ll e^{-c|\gamma|}\,(1+|\gamma|)^k,
  \qquad
  |b_\rho|^2 \ll e^{-2c|\gamma|}\,(1+|\gamma|)^{2k}.
\]
\end{lemma}

\begin{proof}
Write $s=\sigma+it$ with $0\le \sigma\le 1$. By Stirling's formula in vertical strips (see, e.g., \cite[\S2.3]{Titchmarsh}),
\[
  \Bigl|\Gamma\!\Bigl(\frac{s}{2}\Bigr)\Bigr|
  \ll (1+|t|)^{\sigma/2-1/2}\exp\!\Bigl(-\frac{\pi}{4}|t|\Bigr),
\]
and trivially $|\pi^{-s/2}|=\pi^{-\sigma/2}\asymp 1$ on $0\le \sigma\le 1$.
Hence
\[
  |G(\sigma+it)|
  = |\pi^{-s/2}|\cdot \Bigl|\Gamma\!\Bigl(\frac{s}{2}\Bigr)\Bigr|\cdot |H(s)|
  \ll (1+|t|)^{A_H+1}\exp\!\Bigl(-\frac{\pi}{4}|t|\Bigr).
\]
For a zero $\rho=\beta+i\gamma$ we have $0<\beta<1$, and for $|\gamma|\ge 1$,
\[
  \Bigl|\frac{\rho-1}{\rho}\Bigr|
  \le 1 + \frac{1}{|\rho|}
  \le 2.
\]
Therefore,
\[
  |b_\rho|
  = \Bigl|\frac{\rho-1}{\rho}\Bigr|\cdot |G(\rho)|
  \ll (1+|\gamma|)^{A_H+1}\exp\!\Bigl(-\frac{\pi}{4}|\gamma|\Bigr),
\]
which implies the claimed bound with $c=\pi/4$ and $k=A_H+1$ (after adjusting the implied constant to cover the finitely many zeros with $|\gamma|<1$).
\end{proof}

\begin{lemma}[Local zero counting]\label{lem:localcount}
Let $N(T)$ be the number of non-trivial zeros $\rho=\beta+i\gamma$ with $0<\beta<1$ and $0<\gamma\le T$, counted with multiplicity.
Then, for $T\ge 3$,
\[
  N(T)=\frac{T}{2\pi}\log\frac{T}{2\pi}-\frac{T}{2\pi}+O(\log T),
\]
and consequently
\[
  N(T+1)-N(T)\ll \log T.
\]
In particular, for every $y\ge 3$ the interval $(y,y+1]$ contains $\ll \log y$ ordinates $\gamma$ (counted with multiplicity). By symmetry of the zeros under $\gamma\mapsto -\gamma$, the same bound holds for $y\le -3$, hence for all $|y|\ge 3$,
\[
  \#\{\rho:\ y<\Im(\rho)\le y+1\}\ \ll\ \log(|y|+2).
\]
\end{lemma}

\begin{proof}
This is the classical Riemann--von Mangoldt formula; see \cite[Ch.\ 9]{Titchmarsh}. Taking differences gives
\[
  N(T+1)-N(T)
  = \frac{1}{2\pi}\log T + O(1),
\]
which implies $N(T+1)-N(T)\ll \log T$ for $T\ge 3$.
\end{proof}

\begin{lemma}[Large sieve for separated exponentials]\label{lem:largesieve}
Let $\lambda_1,\dots,\lambda_N\in\mathbb{R}$ satisfy $|\lambda_m-\lambda_n|\ge \delta>0$ for all $m\neq n$. Then for any complex coefficients $a_1,\dots,a_N$ and any $L>0$,
\[
  \int_0^L \Bigl|\sum_{n=1}^N a_n e^{i\lambda_n u}\Bigr|^2\,du
  \le \Bigl(L+\frac{2\pi}{\delta}\Bigr)\sum_{n=1}^N |a_n|^2.
\]
\end{lemma}
\begin{proof}
This is \cite[Theorem~1.3]{MontgomeryLNM227}; see also \cite[Ch.\ 7]{MontgomeryVaughan}.
\end{proof}

\section{Theorem A: Dyadic \texorpdfstring{$L^2$}{L2} tail bound}
\begin{theorem}[Dyadic $L^2$ tail bound]\label{thm:A}
Assume Hypotheses \textbf{(H1)}--\textbf{(H3)}. Then there exist constants $c>0$, $k\ge 0$ and $T_0\ge 3$ such that, for all $T\ge T_0$,
\begin{equation}\label{eq:thmA-exp}
  \sum_{\substack{\rho \\ T < |\gamma| \le 2T}} |b_\rho|^2
  \;\ll\; T^{2k+1}(\log T)\,e^{-2cT}.
\end{equation}
In particular, there exist constants $C>0$ and $A>0$ such that, for all $T\ge T_0$,
\[
  \sum_{\substack{\rho \\ T < |\gamma| \le 2T}} |b_\rho|^2 \;\le\; C\,T(\log T)^A.
\]
\end{theorem}

\begin{proof}
By Lemma~\ref{lem:decay}, there exist $c>0$ and $k\ge 0$ such that
\[
  |b_\rho|^2 \ll e^{-2c|\gamma|}\,(1+|\gamma|)^{2k}
\]
for every non-trivial zero $\rho=\beta+i\gamma$.
For $T<|\gamma|\le 2T$ we therefore have the uniform bound
\[
  |b_\rho|^2 \ll e^{-2cT}\,(1+2T)^{2k} \ll e^{-2cT}\,T^{2k}.
\]
Summing over the dyadic band and using the Riemann--von Mangoldt bound
$N(2T)-N(T)\ll T\log T$ (which follows from Lemma~\ref{lem:localcount}) gives
\[
  \sum_{\substack{\rho\\ T<|\gamma|\le 2T}} |b_\rho|^2
  \ll (N(2T)-N(T))\, e^{-2cT}\,T^{2k}
  \ll T^{2k+1}(\log T)\,e^{-2cT}.
\]
In particular the dyadic $L^2$ tail is exponentially small in $T$.
Since $e^{2cT}$ dominates any fixed power of $T$ and $\log T$, there exist constants
$A>0$ and $T_0>0$ such that for all $T\ge T_0$,
\[
  T^{2k}(\log T)\,e^{-2cT} \le (\log T)^A,
\]
and hence
\[
  \sum_{\substack{\rho\\ T<|\gamma|\le 2T}} |b_\rho|^2 \le C\,T(\log T)^A
\]
for a suitable constant $C>0$.
\end{proof}
\begin{remark}
The proof in fact yields an exponentially small dyadic tail in $T$; we record only a polynomial bound because it suffices for applications.
\end{remark}

\section{Theorem B: A large-sieve inequality for zero ordinates}
Fix $L>0$ and a dyadic band $T<|\gamma|\le 2T$ with $T$ large. Define the spectral sum
\[
  F(u) = \sum_{T<|\gamma|\le 2T} b_\rho\, e^{i\gamma u}.
\]

\begin{theorem}[Large-sieve mean-square inequality for zero ordinates]\label{thm:B}
Assume Hypothesis \textbf{(H3)}. Let $T\ge 3$ and $L\ge 1$. For the dyadic band $T<|\gamma|\le 2T$ define
\[
  F(u) = \sum_{\substack{\rho\\ T<|\gamma|\le 2T}} b_\rho\, e^{i\gamma u}.
\]
Then
\begin{equation}\label{eq:thmB}
  \frac{1}{L} \int_0^L |F(u)|^2 \, du
    \;\ll\; (\log T)\sum_{\substack{\rho\\ T<|\gamma|\le 2T}} |b_\rho|^2,
\end{equation}
with an absolute implied constant. In particular, \eqref{eq:thmB} holds in any ``admissible range'' $1\ll L\ll T^\alpha$ with $0<\alpha<1$.
\end{theorem}

\begin{lemma}[Decomposition into $1$-separated subfamilies]\label{lem:partition}
Let $\Lambda$ be a finite multiset of real numbers with the property that every interval of length $1$
contains at most $M$ elements of $\Lambda$ (counted with multiplicity). Then $\Lambda$ can be written
as a disjoint union
\[
  \Lambda=\Lambda_1\cup\cdots\cup \Lambda_M
\]
such that each $\Lambda_j$ is $1$-separated, i.e.\ $|\lambda-\lambda'|\ge 1$ for all distinct $\lambda,\lambda'\in\Lambda_j$.
\end{lemma}

\begin{proof}
Order the elements of $\Lambda$ (with multiplicity) as $\lambda_1\le \lambda_2\le \cdots\le \lambda_N$.
We construct the subfamilies $\Lambda_1,\dots,\Lambda_M$ greedily.

Assume $\lambda_1,\dots,\lambda_{n-1}$ have been assigned so that each $\Lambda_j$ is $1$-separated.
Consider $\lambda_n$. For each $j$, the $1$-separation implies that $\Lambda_j$ contains \emph{at most one}
element in the interval $(\lambda_n-1,\lambda_n]$.
Hence, if $\lambda_n$ cannot be placed into $\Lambda_j$ (i.e.\ would violate $1$-separation), then $\Lambda_j$
already contains an element of $(\lambda_n-1,\lambda_n]$.

But $(\lambda_n-1,\lambda_n]$ contains at most $M$ elements of $\Lambda$ by hypothesis, so at most $M-1$
of the sets $\Lambda_j$ can be ``blocked'' in this way by previously placed elements.
Therefore there exists at least one index $j$ for which $\Lambda_j$ contains no element of $(\lambda_n-1,\lambda_n]$,
and we may assign $\lambda_n$ to that $\Lambda_j$ while preserving $1$-separation.

Proceeding inductively assigns all elements to $\Lambda_1,\dots,\Lambda_M$, as required.
\end{proof}

\begin{proof}
Let $\Gamma(T)$ denote the multiset of ordinates $\gamma\in\mathbb{R}$ (counted with multiplicity) of zeros
$\rho=\beta+i\gamma$ with $T<|\gamma|\le 2T$.
By Lemma~\ref{lem:localcount}, every interval of length $1$ in the ranges $[T,2T]$ and $[-2T,-T]$
contains $\ll \log T$ ordinates. Thus there exists an absolute constant $C_0>0$ such that
\[
  M:=\sup_{y\in\mathbb{R}}\#\bigl(\Gamma(T)\cap (y,y+1]\bigr) \le C_0\log T
\]
for all sufficiently large $T$.

Apply Lemma~\ref{lem:partition} to $\Gamma(T)$ to obtain a disjoint union
$\Gamma(T)=\Gamma_1\cup\cdots\cup \Gamma_M$, where each $\Gamma_j$ is $1$-separated.
Write accordingly
\[
  F(u)=\sum_{j=1}^M F_j(u),
  \qquad
  F_j(u):=\sum_{\substack{\rho\\ \gamma\in\Gamma_j}} b_\rho\,e^{i\gamma u}.
\]
By Cauchy--Schwarz,
\[
  |F(u)|^2 = \Bigl|\sum_{j=1}^M F_j(u)\Bigr|^2 \le M \sum_{j=1}^M |F_j(u)|^2,
\]
and hence
\begin{equation}\label{eq:splitCS}
  \int_0^L |F(u)|^2\,du \le M \sum_{j=1}^M \int_0^L |F_j(u)|^2\,du.
\end{equation}

For each fixed $j$, the ordinates in $\Gamma_j$ are $1$-separated, so Lemma~\ref{lem:largesieve} with $\delta=1$ gives
\[
  \int_0^L |F_j(u)|^2\,du \le (L+2\pi)\sum_{\substack{\rho\\ \gamma\in\Gamma_j}} |b_\rho|^2.
\]
Substituting into \eqref{eq:splitCS} and summing over $j$ yields
\[
  \int_0^L |F(u)|^2\,du
  \le M(L+2\pi)\sum_{T<|\gamma|\le 2T} |b_\rho|^2.
\]
Dividing by $L$ and using $L\ge 1$ gives
\[
  \frac{1}{L}\int_0^L |F(u)|^2\,du
  \le C\,(\log T)\sum_{T<|\gamma|\le 2T} |b_\rho|^2,
\]
for a suitable absolute constant (absorbing the factor $(1+2\pi/L)$). This is exactly \eqref{eq:thmB}.
\end{proof}

\begin{corollary}[Max-bound version]\label{cor:max}
Under the assumptions of Theorem~\ref{thm:B}, one has
\[
  \frac{1}{L} \int_0^L |F(u)|^2 \, du
    \;\ll\; T \,(\log T)^{2} \cdot \max_{T<|\gamma|\le 2T} |b_\rho|^2.
\]
\end{corollary}
\begin{proof}
This follows immediately from Theorem \ref{thm:B} and the trivial bound $\sum |b_\rho|^2 \le (N(2T)-N(T)) \max |b_\rho|^2 \ll T (\log T) \max |b_\rho|^2$.
\end{proof}

\section{Theorem C: A uniform RMS tail bound}
Let $S\ge 3$ and $L>0$, and fix a parameter $B>0$. Define the tail observable
\[
  \widetilde F_{\mathrm{tail}}(s) \;:=\; s^{-2} \sum_{\substack{|\gamma|>S^B}} b_\rho\, e^{i\gamma s},
\]
and the windowed root-mean-square quantity
\[
  K_{\mathrm{tail}}(S,L)^2 \;:=\; \frac{1}{L} \int_{S}^{S+L} \bigl|\widetilde F_{\mathrm{tail}}(s)\bigr|^2 \, ds.
\]

\begin{theorem}[Uniform RMS tail bound]\label{thm:C}
Assume Hypotheses \textbf{(H1)}--\textbf{(H3)} and fix $B>0$. Then there exist constants $c'>0$ and $S_0\ge 3$ such that, for all $S\ge S_0$ and all $L>0$,
\begin{equation}\label{eq:thmC-exp}
  K_{\mathrm{tail}}(S,L)^2 \;\ll\; S^{-4}\exp\!\bigl(-2c' S^{B}\bigr).
\end{equation}
In particular, $K_{\mathrm{tail}}(S,L)^2\le M_{\mathrm{tail}}^2$ uniformly in $S$ and $L$, for a suitable constant $M_{\mathrm{tail}}>0$.
\end{theorem}

\begin{proof}
By Lemma~\ref{lem:decay}, there exist $c>0$ and $k\ge 0$ such that
\[
  |b_\rho| \ll e^{-c|\gamma|}(1+|\gamma|)^{k}.
\]
We first show that the tail sum $\sum_{|\gamma|>S^B}|b_\rho|$ is exponentially small in $S^B$.
Decompose the region $|\gamma|>S^B$ into dyadic bands $(2^m S^B,2^{m+1}S^B]$, $m\ge 0$. On each band we have
$|b_\rho|\ll e^{-c2^mS^B}(2^mS^B)^k$, and by Lemma~\ref{lem:localcount} (equivalently, by Riemann--von Mangoldt),
\[
  \#\{\rho:\ 2^mS^B<|\gamma|\le 2^{m+1}S^B\}
  \;=\; O\!\bigl(2^mS^B\log(2^mS^B)\bigr).
\]
Hence
\[
  \sum_{\substack{\rho\\ 2^mS^B<|\gamma|\le 2^{m+1}S^B}} |b_\rho|
  \ll (2^mS^B)^{k+1}\log(2^mS^B)\,e^{-c2^mS^B}.
\]
Summing over $m\ge 0$ gives
\[
  \sum_{\substack{\rho\\ |\gamma|>S^B}} |b_\rho|
  \ll \sum_{m=0}^\infty (2^mS^B)^{k+1}\log(2^mS^B)\,e^{-c2^mS^B}
  \ll \exp(-c' S^B)
\]
for some $c'>0$ and all $S\ge S_0$, since the exponential factor decays super-geometrically in $m$.

Now for $s\in[S,S+L]$ we have $|s^{-2}|\le S^{-2}$, and by absolute convergence,
\[
  |\widetilde F_{\mathrm{tail}}(s)|
  \le S^{-2}\sum_{|\gamma|>S^B}|b_\rho|
  \ll S^{-2}\exp(-c'S^B).
\]
Therefore
\[
  K_{\mathrm{tail}}(S,L)^2
  =\frac1L\int_S^{S+L}|\widetilde F_{\mathrm{tail}}(s)|^2\,ds
  \ll \frac1L\int_S^{S+L} S^{-4}e^{-2c'S^B}\,ds
  \ll S^{-4}e^{-2c'S^B},
\]
which is \eqref{eq:thmC-exp}.
\end{proof}

\begin{remark}
Because the coefficients $b_\rho$ decay exponentially in $|\gamma|$ under \textbf{(H1)}--\textbf{(H2)}, the tail bound above is essentially pointwise and does not require any mean-square input beyond absolute convergence. One may also derive comparable bounds by combining dyadic decompositions with Theorems~\ref{thm:A} and \ref{thm:B}, but the direct argument above is simpler.
\end{remark}

\section{Remarks}

\begin{itemize}
  \item The proofs use only classical inputs: Stirling's formula for $\Gamma$, the Riemann--von Mangoldt formula (including unit-interval zero counts), and the large sieve for separated exponentials (Lemma~\ref{lem:largesieve}).
  \item Under \textbf{(H1)}--\textbf{(H2)}, the dyadic tails in Theorem~\ref{thm:A} are in fact exponentially small in the dyadic parameter.
  \item The large-sieve inequality in Theorem~\ref{thm:B} does not assume any uniform lower bound on gaps between ordinates; it uses only the local counting information encoded in \textbf{(H3)}.
  \item Variants with different test functions or normalizations of $b_\rho$ can be handled similarly by modifying $H(s)$ and verifying \textbf{(H1)}--\textbf{(H2)} for the resulting factor.
\end{itemize}

\begin{thebibliography}{9}

\bibitem{Titchmarsh}
E.\,C.~Titchmarsh,
\emph{The Theory of the Riemann Zeta-Function},
2nd ed., revised by D.\,R.~Heath-Brown,
Oxford University Press, 1986.

\bibitem{MontgomeryVaughan}
H.\,L.~Montgomery and R.\,C.~Vaughan,
\emph{Multiplicative Number Theory I.\ Classical Theory},
Cambridge University Press, 2007.

\bibitem{MontgomeryLNM227}
H.\,L.~Montgomery,
\emph{Topics in Multiplicative Number Theory},
Lecture Notes in Mathematics 227,
Springer, 1971.

\end{thebibliography}

\end{document}
